\section{Classical PBFT and Its Performance Characteristics}

Classical PBFT~\cite{castro1999practical} established the design template for practical Byzantine agreement and remains the baseline against which later protocols are measured. It operates in the partially synchronous model, assuming that message delays become bounded after an unknown Global Stabilization Time (GST), and tolerates $f$ Byzantine failures among $n$ replicas subject to the safety requirement $n > 3f$. Within these assumptions the protocol delivers deterministic termination, rather than the probabilistic guarantees to which the FLP impossibility result confines fully asynchronous designs. The bound is not an artifact of the protocol: with $n \leq 3f$, Byzantine nodes can partition the network into groups observing contradictory but internally consistent states, so the constraint applies to every deterministic Byzantine consensus protocol regardless of timing assumptions or communication pattern.

\subsection{The Three-Phase Commit Protocol}

Agreement proceeds through three phases: pre-prepare, prepare, and commit. In the pre-prepare phase, the primary of the current view assigns a sequence number to a client request and broadcasts the resulting proposal to the replica set, fixing a candidate order. The prepare phase has replicas echo the proposal to one another; a replica becomes \emph{prepared} once it observes matching prepare messages from a quorum, which certifies that no conflicting proposal can accumulate a competing quorum within the same view. The commit phase performs a second all-to-all exchange, establishing that the prepared decision is held by enough correct replicas to survive a change of primary. The division of labour between the two exchanges is the essential structural point: the prepare phase secures agreement \emph{within} a view, while the commit phase makes that agreement durable \emph{across} views. Safety thus reduces to quorum intersection under $n > 3f$, and any later protocol that collapses or pipelines these exchanges must recover the cross-view guarantee by other means.

\subsection{View Change and Liveness}

Because the primary may itself be Byzantine, progress cannot depend on it. Replicas arm timers on pending requests and, on expiry, abandon view $v$ for view $v+1$, transmitting to the incoming primary the certificates witnessing their highest prepared decisions. The incoming primary must reconstruct a safe prefix from a quorum of such reports before proposing again. This mechanism converts a safety-preserving core into a live protocol once the network stabilises after GST, but it is also the least uniform component of the design: view-change messages carry accumulated certificate state rather than a single fixed-size proposal, so their cost is not represented by the steady-state message count and grows with both replica count and the depth of uncommitted state.

\subsection{Communication Complexity as a Baseline Metric}

The steady-state cost is $\mathcal{O}(n^2)$ messages per consensus instance, dominated by the two all-to-all phases, with a normal-case latency of three communication rounds after the client request reaches the primary. This quadratic term is the canonical baseline figure for the protocol family; subsequent work re-expresses it as an authenticator complexity so that signature generation, transmission, and verification costs become visible alongside raw message counts~\cite{yin2019hotstuff}.

\subsection{Scalability Bottlenecks}

Three pressure points follow from the analysis. First, quadratic message growth makes bandwidth, not computation, the binding constraint as $n$ increases, which limits practical deployments to modest replica counts. Second, load is leader-centric: the primary originates every proposal and absorbs the view-change reports, so its capacity bounds system throughput while the remaining replicas are underutilised. Third, the view change is non-uniform and unpipelined, making recovery cost qualitatively different from normal-case cost and complicating both performance prediction and formal analysis. These bottlenecks define the optimisation targets pursued by Tendermint~\cite{buchman2016tendermint}, HotStuff~\cite{yin2019hotstuff}, and Casper~\cite{buterin2020combining}, examined in the sections that follow; empirical throughput and latency figures from production deployments are deferred to the evaluation discussion.